
\documentclass[11pt,a4paper,spanish]{book}
\usepackage{estilo_unir-1}
\hypersetup{
	colorlinks   = true, %Colours links instead of ugly boxes
	urlcolor     = blue, %Colour for external hyperlinks
	linkcolor    = black, %Colour of internal links
	citecolor   = red %Colour of citations
}
\usepackage[nottoc]{tocbibind}

%---------------------------
%título de la asignatura, trabajo y autor
%---------------------------
\subject{Trabajo de Fin de Master}
\title{Marco Híbrido GNN-HVA para la Caracterización de Fases Cuánticas en Hardware NISQ: Aceleración VQE mediante Redes Neuronales de Grafos y Ansätze Variacionales Superficiales}
\author{Franco Raineri}
\profesor{Andres Navas Caliz}
\date{15 de Abril del
 2026}

%---------------------------
%marges
%---------------------------
%\usepackage[margin=1.9cm]{geometry}
%---------------------------
%---------------------------
%---------------------------
%---------------------------
\begin{document}
\renewcommand{\listfigurename}{Índice de figuras}
\renewcommand{\listtablename}{Índice de tablas}
\renewcommand{\contentsname}{Índice de contenidos}
\renewcommand{\figurename}{Figura}
\renewcommand{\tablename}{Tabla} 

\maketitle
\frontmatter
\tableofcontents
\listoffigures

\listoftables


\chapter{Resumen}


\vspace{0.5cm}
{\bf Palabras clave:} 

\chapter{Abstract}

\vspace{0.5cm}
{\bf Keywords:}

\mainmatter

\chapter{Introducción}

En el ámbito científico actual, la resolución del Problema de los Muchos Cuerpos (Quantum Many-Body Problem) representa uno de los desafíos centrales y más estimulantes de la física de la materia condensada. Cuando un gran número de partículas interactúan a nivel cuántico, el sistema no puede describirse simplemente como la suma de sus partes aisladas; en su lugar, emergen comportamientos colectivos y leyes físicas fundamentalmente nuevas \citep{anderson1972}. El foco de esta investigación radica en una de las manifestaciones más relevantes de este fenómeno: la caracterización de transiciones de fase cuánticas, donde el sistema experimenta cambios cualitativos en su estado fundamental al variar un parámetro externo, sin intervención de fluctuaciones térmicas \citep{dutta2015}.

Durante décadas, la física comprendió las fases de la materia a través de la Teoría de Ruptura de Simetría de Landau, donde las transiciones de fase se describen mediante la aparición de un parámetro de orden local medible. Sin embargo, a temperaturas cercanas al cero absoluto, donde las fluctuaciones térmicas desaparecen, emergen fases de la materia que desafían este paradigma. Las transiciones de fase cuánticas (QPT) están gobernadas por fluctuaciones puramente cuánticas y presentan fenómenos de entrelazamiento de largo alcance que hacen intratable su simulación clásica. En el caso más exótico, las fases topológicas ---como los Líquidos de Espín Cuántico (QSL)--- carecen por completo de un orden local, codificando sus propiedades en correlaciones cuánticas globales \citep{savary2016}.

El problema fundamental es que simular el comportamiento de estos sistemas altamente correlacionados resulta matemáticamente intratable para los simuladores clásicos convencionales. El espacio de Hilbert que describe estos sistemas crece de manera exponencial con el número de partículas, un fenómeno conocido como la ``Maldición de la Dimensionalidad''. Además, los métodos estadísticos diseñados para eludir esta barrera (como el Monte Carlo Cuántico) fracasan ante el ``Problema del Signo'', haciendo que la simulación clásica exacta sea inviable para sistemas frustrados grandes \citep{troyer2005}.

El presente Trabajo Fin de Máster (TFM) propone el diseño, justificación e implementación de un marco de simulación híbrido clásico-cuántico. Este marco integra técnicas de Inteligencia Artificial clásica ---específicamente Redes Neuronales de Grafos de Paso de Mensajes (MPNN)--- con algoritmos cuánticos variacionales ejecutados en circuitos superficiales, diseñados para operar dentro de las limitaciones del hardware cuántico ruidoso de escala intermedia (NISQ) \citep{preskill2018}. El objetivo central es demostrar que la predicción clásica de parámetros variacionales, combinada con un Ansatz Variacional Hamiltoniano (HVA) de profundidad estrictamente acotada, permite clasificar correctamente fases cuánticas mediante observables locales con un costo cuántico cercano a cero ---reduciendo de cientos de evaluaciones en hardware a 0--2 iteraciones de refinamiento.

Como sistema de validación se emplea el Modelo de Ising en Campo Transversal (TFIM) unidimensional, un sistema paradigmático que exhibe una transición de fase cuántica entre las fases paramagnética y ferromagnética. Este modelo permite verificar rigurosamente el pipeline en un régimen donde la solución exacta es conocida, estableciendo la metodología que escala sin modificaciones de código a sistemas de mayor tamaño donde los métodos clásicos fallan \citep{tripathi2026}.

A continuación, se detalla la justificación del problema específico a tratar, el planteamiento riguroso de la solución algorítmica propuesta y la estructura en la que se organizan de manera lógica los capítulos posteriores de esta memoria investigativa.

\section{Motivación}

El problema central que este Trabajo pretende abordar es la incapacidad de los algoritmos cuánticos variacionales convencionales para operar eficientemente en el hardware NISQ actual, específicamente en el contexto de la caracterización de fases cuánticas de la materia. La simulación cuántica fue propuesta originalmente por \citet{feynman1982} como la solución natural al crecimiento exponencial del espacio de Hilbert, y los algoritmos variacionales como VQE \citep{peruzzo2014} representan la aproximación más prometedora para dispositivos de escala intermedia.

Sin embargo, tres barreras fundamentales limitan la aplicabilidad práctica de VQE en hardware real:

\begin{enumerate}
    \item \textbf{Barren Plateaus:} La optimización de circuitos cuánticos parametrizados sufre la desaparición exponencial de gradientes cuando se emplean funciones de costo globales o circuitos profundos aleatorios \citep{mcclean2018, cerezo2021}. Este fenómeno hace que el optimizador clásico no encuentre dirección de descenso válida, imposibilitando el entrenamiento.
    
    \item \textbf{Truncamiento por Ruido:} Hallazgos recientes han demostrado rigurosamente que el ruido no unital inherente en el hardware actual (relajación térmica, amortiguamiento de amplitud) trunca la profundidad efectiva de cualquier circuito cuántico a $O(\log n)$ capas \citep{mele2026}. En la práctica, toda información codificada más allá de esta profundidad se pierde irremediablemente, independientemente de la estrategia de optimización empleada.
    
    \item \textbf{Costo de Evaluación:} Un barrido VQE convencional requiere cientos a miles de evaluaciones de circuito por punto del espacio de parámetros. En hardware cuántico, cada evaluación consume tiempo de coherencia y presupuesto de shots, haciendo prohibitivo el mapeo completo de un diagrama de fases.
\end{enumerate}

La motivación de este trabajo surge de la convergencia de tres resultados teóricos recientes que, combinados, abren una ventana de operación viable:

\begin{itemize}
    \item Los circuitos superficiales con funciones de costo locales están \emph{probadamente} libres de barren plateaus \citep{cerezo2021, mele2026}.
    \item El Ansatz Variacional Hamiltoniano (HVA) respeta las simetrías del Hamiltoniano, maximizando la expresividad por parámetro en circuitos de baja profundidad \citep{wiersema2020}.
    \item Las Redes Neuronales de Grafos pueden aprender eficientemente propiedades del estado fundamental dentro de una fase, generalizando a nuevas instancias del Hamiltoniano \citep{huang2022}.
\end{itemize}

La relevancia a largo plazo de este enfoque radica en su extensibilidad hacia sistemas frustrados bidimensionales ---como las redes de Kagomé, donde se han demostrado recientemente simulaciones de 103 sitios en hardware IBM \citep{ahsan2025}--- y hacia fases topológicas donde las excitaciones fraccionalizadas (anyones) posibilitan memorias cuánticas tolerantes a fallos \citep{savary2016}.

\section{Planteamiento del problema}

Para sortear simultáneamente las tres barreras identificadas, este Trabajo propone abandonar tanto el uso de ansätze genéricos tipo Hardware-Efficient (HEA) ---cuya inferioridad frente al HVA ha sido confirmada independientemente para el TFIM en 1D, 2D y 3D hasta 27 espines \citep{tripathi2026}--- como la optimización directa en hardware sin inicialización informada.

En su lugar, se propone el desarrollo de una \textbf{Arquitectura Predictiva Híbrida (GNN-HVA)}, implementada como un pipeline de cuatro fases:

\begin{enumerate}
    \item \textbf{Fase 1 --- Ground Truth Clásico:} Se genera un conjunto de datos de referencia mediante diagonalización exacta y DMRG \citep{schollwock2011}, obteniendo estados fundamentales, energías y observables locales del TFIM para un barrido denso del parámetro de campo transversal $h$.
    
    \item \textbf{Fase 2 --- Optimización VQE con Warm-Start:} Se construye un Ansatz Variacional Hamiltoniano (HVA) de profundidad $p \leq 2$ capas \citep{wiersema2020} y se optimizan sus parámetros mediante un barrido descendente $h = 2 \to 0$, donde cada punto hereda los parámetros óptimos del anterior como semilla inicial. Esta estrategia de warm-start reduce el costo de optimización en un factor $10$--$50\times$ respecto a la inicialización aleatoria \citep{puig2025, mele2022}.
    
    \item \textbf{Fase 3 --- Predictor MPNN:} Se entrena una Red Neuronal de Grafos de Paso de Mensajes (MPNN) basada en GINConv \citep{xu2019} para aprender el mapeo directo desde la topología del Hamiltoniano (representado como grafo) hacia los parámetros óptimos $\theta^*$ del circuito HVA. Una vez entrenada, la red predice parámetros en una única pasada forward clásica, eliminando la necesidad de optimización cuántica iterativa.
    
    \item \textbf{Fase 4 --- Despliegue en Hardware:} Los parámetros predichos por la MPNN se inyectan como semilla en el circuito HVA, que se ejecuta en hardware cuántico real (IBM Torino/Heron) con mitigación de errores: Zero Noise Extrapolation (ZNE) \citep{uvarov2024}, Dynamical Decoupling \citep{pokharel2025} y Twirled Readout Error eXtinction (TREX). El criterio de éxito es $\Delta E / \text{gap} < 5\%$ y clasificación correcta de fase mediante observables locales $\langle X_i \rangle$ y $\langle Z_i Z_{i+1} \rangle$.
\end{enumerate}

El objetivo general es demostrar que este pipeline híbrido permite clasificar correctamente fases cuánticas en hardware NISQ con un costo cuántico de 0--2 iteraciones de refinamiento adaptativo, frente a las 100--1000 evaluaciones requeridas por VQE convencional. Esto constituye una reducción de $50$--$500\times$ en tiempo de QPU, validando la viabilidad del paradigma de inicialización inteligente para algoritmos variacionales \citep{zhang2025, miao2024}.

\section{Estructura del Trabajo}

Para abordar sistemáticamente los objetivos descritos y proporcionar una demostración rigurosa, reproducible y secuencial del marco híbrido propuesto, el resto de la presente memoria se ha estructurado en cinco capítulos principales, organizados de la siguiente manera:

\begin{itemize}
    \item \textbf{Capítulo 2: Marco Teórico y Estado del Arte.} Este capítulo sienta las bases físicas y matemáticas de la investigación. Se introduce formalmente el Modelo de Ising en Campo Transversal (TFIM) y su transición de fase cuántica, el formalismo del Eigensolver Cuántico Variacional (VQE) y la construcción del Ansatz Variacional Hamiltoniano (HVA). Se realiza un análisis crítico del estado del arte, examinando las barreras teóricas de los \emph{barren plateaus} \citep{mcclean2018, cerezo2021}, el truncamiento de profundidad por ruido no unital \citep{mele2026}, y las aproximaciones recientes de inicialización mediante aprendizaje automático \citep{miao2024, zhang2025, zou2026}. Se introduce además la teoría de Redes Neuronales de Grafos y su aplicación a sistemas de espines \citep{xu2019, gilmer2017, kochkov2021}.
    
    \item \textbf{Capítulo 3: Metodología.} Se describe la arquitectura completa del pipeline de cuatro fases: generación de ground truth clásico (diagonalización exacta + DMRG), optimización VQE con warm-start descendente, entrenamiento del predictor MPNN (GINConv, filtro de fidelidad $\geq 0.93$, validación física), y despliegue en hardware con mitigación de errores (ZNE inhomogéneo, DD, TREX). Se detalla el marco de validación de seis métricas con ordenamiento por prioridad y los criterios de éxito para clasificación de fases.
    
    \item \textbf{Capítulo 4: Resultados.} Se presentan los resultados experimentales organizados por escala del sistema. Para $N=6$ (cadena 1D): optimización de hiperparámetros sobre 40+ experimentos y 14 configuraciones, alcanzando checklist 4/4 en $h \geq 1.25$. Para $N=10$: validación de escalabilidad con configuración óptima (MPNN $h_{\text{dim}}=128$, seed 43, patience 500), alcanzando $\Delta E/\text{gap} = 4.79\% \pm 0.50\%$ en $h=1.4$. Se incluye el análisis del límite físico en $h=1.25$ como frontera de expresividad del HVA $p=2$ ---confirmado independientemente por \citet{tripathi2026}--- y los resultados de despliegue en hardware con mitigación de errores.
    
    \item \textbf{Capítulo 5: Discusión.} Se analiza críticamente la validación del pipeline, comparando con trabajos contemporáneos (Qracle \citep{zhang2025}, NN-VQE \citep{miao2024}, Flow-VQE \citep{zou2026}). Se discuten las limitaciones y fronteras físicas del enfoque, el argumento de utilidad cuántica, y las direcciones futuras incluyendo la extensión a redes de Kagomé y la incorporación de técnicas generativas.
    
    \item \textbf{Capítulo 6: Conclusiones y Trabajo Futuro.} El capítulo final sintetiza las contribuciones de la investigación, evalúa el grado de cumplimiento de los objetivos, y propone líneas de investigación futuras para la extensión del pipeline a sistemas bidimensionales frustrados y hardware cuántico tolerante a fallos.
\end{itemize}


%---------------------------
% CAPÍTULO 2: MARCO TEÓRICO Y ESTADO DEL ARTE
%---------------------------

\chapter{Marco Teórico y Estado del Arte}

Este capítulo presenta los fundamentos teóricos y el estado del conocimiento en las tres disciplinas que convergen en este trabajo: la física de transiciones de fase cuánticas, los algoritmos cuánticos variacionales y las redes neuronales de grafos. El objetivo es establecer el marco conceptual necesario para comprender la arquitectura híbrida propuesta, identificar las limitaciones del estado del arte actual y justificar las decisiones de diseño adoptadas.

\section{El Problema de los Muchos Cuerpos Cuántico}

El Problema de los Muchos Cuerpos Cuántico (Quantum Many-Body Problem) constituye uno de los desafíos computacionales fundamentales de la física moderna. Dado un sistema de $N$ partículas cuánticas interactuantes, su estado se describe mediante un vector en un espacio de Hilbert cuya dimensión crece exponencialmente como $2^N$ para espines-1/2. Para $N = 40$ partículas, el espacio de estados alcanza $\sim 10^{12}$ dimensiones, superando la capacidad de almacenamiento de cualquier supercomputador clásico existente \citep{feynman1982}.

Este crecimiento exponencial no es meramente un problema de recursos computacionales, sino una barrera fundamental. \citet{troyer2005} demostraron que la simulación general de sistemas cuánticos fermionicos mediante Monte Carlo es NP-hard debido al Problema del Signo: las amplitudes de probabilidad pueden ser negativas, causando cancelaciones catastróficas que impiden la convergencia estadística. Este resultado establece que, salvo $P = NP$, no existe algoritmo clásico eficiente para simular la dinámica general de estos sistemas.

\subsection{Transiciones de Fase Cuánticas}

A diferencia de las transiciones de fase clásicas (gobernadas por fluctuaciones térmicas), las transiciones de fase cuánticas (QPT) ocurren a temperatura cero y están impulsadas por fluctuaciones puramente cuánticas al variar un parámetro del Hamiltoniano \citep{dutta2015}. En el punto crítico cuántico, las correlaciones divergen y el sistema exhibe entrelazamiento de largo alcance, haciendo que los métodos de campo medio y las aproximaciones perturbativas fallen simultáneamente.

La clasificación de fases cuánticas se realiza mediante parámetros de orden: cantidades observables que toman valores distintos en cada fase y se anulan (o divergen) en el punto de transición. Para las transiciones de ruptura de simetría ---como la del Modelo de Ising--- estos parámetros de orden son observables locales, lo cual resulta crucial para su medición en hardware cuántico ruidoso.

\subsection{El Modelo de Ising en Campo Transversal (TFIM)}

El TFIM es el modelo paradigmático para el estudio de transiciones de fase cuánticas. Su Hamiltoniano para una cadena unidimensional de $N$ espines con condiciones de contorno abiertas se escribe como:

\begin{equation}
H = -J \sum_{i=1}^{N-1} Z_i Z_{i+1} - h \sum_{i=1}^{N} X_i
\label{eq:tfim}
\end{equation}

\noindent donde $J > 0$ es la constante de acoplamiento entre espines vecinos, $h$ es la intensidad del campo magnético transversal, y $Z_i$, $X_i$ son las matrices de Pauli actuando sobre el sitio $i$.

El modelo exhibe dos fases bien diferenciadas \citep{dutta2015}:

\begin{itemize}
    \item \textbf{Fase ferromagnética} ($h/J \ll 1$): Los espines se alinean a lo largo del eje $z$, rompiendo espontáneamente la simetría $\mathbb{Z}_2$. El parámetro de orden es la magnetización $\langle Z_i \rangle \neq 0$.
    \item \textbf{Fase paramagnética} ($h/J \gg 1$): El campo transversal domina y los espines se alinean con el campo externo. La magnetización transversal $\langle X_i \rangle \to 1$ y las correlaciones $\langle Z_i Z_{i+1} \rangle \to 0$.
\end{itemize}

En el límite termodinámico ($N \to \infty$), la transición de fase ocurre en el punto crítico exacto $h_c/J = 1.0$. Para sistemas finitos, la transición se suaviza en una región de crossover cuyo ancho escala como $\sim N^{-1}$ \citep{dutta2015}. Esta propiedad de tamaño finito es relevante para nuestro trabajo: a $N = 6$ y $N = 10$, la región crítica se extiende aproximadamente entre $h \in [0.8, 1.4]$.

La elección del TFIM como sistema de validación se justifica por tres razones fundamentales:

\begin{enumerate}
    \item \textbf{Solución exacta conocida:} Permite verificación rigurosa del pipeline completo contra resultados analíticos.
    \item \textbf{Mapeo isomórfico a qubits:} Cada espín-1/2 corresponde directamente a un qubit, sin overhead de codificación (a diferencia de sistemas fermiónicos que requieren transformaciones de Jordan-Wigner).
    \item \textbf{Observables locales suficientes:} La transición se caracteriza completamente mediante $\langle X_i \rangle$ y $\langle Z_i Z_{i+1} \rangle$, que son medibles con presupuestos polinomiales de shots en hardware cuántico.
\end{enumerate}

\section{Algoritmos Cuánticos Variacionales}

\subsection{El Eigensolver Cuántico Variacional (VQE)}

El VQE, propuesto por \citet{peruzzo2014}, es un algoritmo híbrido clásico-cuántico diseñado para encontrar el estado fundamental de un Hamiltoniano $H$. Su principio de funcionamiento se basa en el principio variacional de la mecánica cuántica:

\begin{equation}
E_0 \leq \langle \psi(\boldsymbol{\theta}) | H | \psi(\boldsymbol{\theta}) \rangle = E(\boldsymbol{\theta})
\label{eq:variational}
\end{equation}

\noindent donde $|\psi(\boldsymbol{\theta})\rangle = U(\boldsymbol{\theta})|0\rangle$ es un estado preparado por un circuito cuántico parametrizado $U(\boldsymbol{\theta})$, y $\boldsymbol{\theta} \in \mathbb{R}^p$ son los parámetros a optimizar. El algoritmo alterna entre: (1) preparación y medición del estado en el procesador cuántico, y (2) actualización de parámetros mediante un optimizador clásico.

La revisión exhaustiva de \citet{tilly2022} establece que la eficacia de VQE depende críticamente de tres factores: la elección del ansatz (expresividad vs. entrenabilidad), la función de costo (local vs. global) y la estrategia de inicialización de parámetros. Nuestro trabajo aborda simultáneamente estos tres factores.

\subsection{El Ansatz Variacional Hamiltoniano (HVA)}

El HVA, introducido por \citet{wiersema2020}, construye el circuito parametrizado directamente a partir de la estructura del Hamiltoniano objetivo. Para el TFIM (Ecuación~\ref{eq:tfim}), el HVA de $p$ capas toma la forma:

\begin{equation}
|\psi(\boldsymbol{\theta})\rangle = \prod_{l=1}^{p} \left[ e^{-i\theta_{l,x} \sum_i X_i} \cdot e^{-i\theta_{l,zz} \sum_{\langle i,j \rangle} Z_i Z_j} \right] |+\rangle^{\otimes N}
\label{eq:hva}
\end{equation}

\noindent donde $|+\rangle^{\otimes N}$ es el estado inicial (producto de autoestados de $X$), que corresponde al estado fundamental exacto en el límite $h \to \infty$. Cada capa $l$ aplica secuencialmente la evolución bajo los términos $ZZ$ y $X$ del Hamiltoniano con ángulos independientes $\theta_{l,zz}$ y $\theta_{l,x}$.

Las ventajas del HVA frente a ansätze genéricos tipo Hardware-Efficient (HEA) son múltiples y han sido confirmadas empíricamente por \citet{tripathi2026} en un estudio comparativo exhaustivo sobre TFIM en 1D, 2D y 3D hasta 27 espines:

\begin{itemize}
    \item \textbf{Respeto de simetrías:} El HVA preserva las simetrías del Hamiltoniano por construcción, restringiendo el espacio de búsqueda al subespacio físicamente relevante.
    \item \textbf{Expresividad por parámetro:} Con solo $2p$ parámetros (frente a $O(Np)$ del HEA), el HVA alcanza fidelidades comparables o superiores.
    \item \textbf{Paisaje de optimización suave:} La transferibilidad de parámetros entre Hamiltonianos cercanos genera un paisaje energético sin mínimos locales espurios \citep{mele2022}.
    \item \textbf{Ausencia de barren plateaus:} Para profundidades acotadas, el HVA exhibe gradientes polinomialmente decrecientes (no exponenciales) \citep{wiersema2020}.
\end{itemize}

\subsection{Barren Plateaus: El Obstáculo Central de los Algoritmos Variacionales}

El fenómeno de barren plateaus (BP) fue identificado por \citet{mcclean2018} como la desaparición exponencial de los gradientes de la función de costo con el número de qubits:

\begin{equation}
\text{Var}\left[\frac{\partial C}{\partial \theta_k}\right] \leq F(n) \in O(2^{-n})
\label{eq:bp}
\end{equation}

Este resultado implica que para $n = 20$ qubits, la señal de gradiente es $\sim 10^{-6}$, indistinguible del ruido estadístico con cualquier número práctico de mediciones. Los BP constituyen la barrera de entrenabilidad más severa para los algoritmos variacionales.

\citet{cerezo2021} refinaron este resultado demostrando que los BP dependen críticamente de la \emph{localidad} de la función de costo. Para funciones de costo compuestas por observables locales (actuando sobre $k$ qubits contiguos), la varianza del gradiente satisface:

\begin{equation}
\text{Var}\left[\frac{\partial C_{\text{local}}}{\partial \theta_k}\right] \geq \Omega\left(\frac{1}{\text{poly}(n)}\right)
\label{eq:local_bp}
\end{equation}

\noindent para circuitos de profundidad $O(\log n)$. Esta cota inferior polinomial garantiza que el optimizador siempre dispone de una dirección de descenso detectable, independientemente del tamaño del sistema. Este resultado teórico fundamenta nuestra elección de observables locales ($\langle X_i \rangle$, $\langle Z_i Z_{i+1} \rangle$) como función de costo y métricas de validación.

\subsection{Truncamiento de Profundidad por Ruido No Unital}

El resultado más reciente y determinante para la arquitectura de este trabajo es el de \citet{mele2026}, publicado en Nature Physics. Los autores demuestran que los canales de ruido no unital ---que incluyen la relajación térmica ($T_1$) y el amortiguamiento de amplitud, dominantes en hardware superconductor--- poseen un punto fijo único (típicamente el estado térmico o el estado maximalmente mezclado). Formalmente, para un circuito de profundidad $d$ con ruido no unital de intensidad $\epsilon$ por compuerta:

\begin{equation}
\|\rho_d - \rho_{\text{fijo}}\|_1 \leq e^{-\Omega(\epsilon \cdot d)}
\label{eq:truncation}
\end{equation}

Cuando $d \gg O(\log n / \epsilon)$, el estado de salida es exponencialmente indistinguible del punto fijo del ruido, independientemente del estado inicial o los parámetros del circuito. Este resultado tiene una consecuencia arquitectónica inmediata: \emph{cualquier circuito cuántico en hardware ruidoso es efectivamente superficial}. No existe beneficio en añadir capas más allá de $O(\log n)$, ya que la información se pierde irremediablemente.

Crucialmente, \citet{mele2026} demuestran también el resultado complementario positivo: en este régimen de circuitos efectivamente superficiales con ruido no unital, las funciones de costo locales \emph{no} sufren barren plateaus. Esto establece la sinergia teórica fundamental de nuestra arquitectura:

\begin{enumerate}
    \item Circuitos superficiales ($p \leq 2$) $\to$ sobreviven al truncamiento por ruido.
    \item Funciones de costo locales $\to$ gradientes detectables (sin BP).
    \item Estructura HVA $\to$ máxima expresividad por parámetro en el régimen superficial.
\end{enumerate}

Esta triple convergencia no es meramente conveniente: es el \emph{único} régimen conocido donde los algoritmos variacionales son simultáneamente entrenables, resilientes al ruido y físicamente significativos en hardware NISQ.

\section{Estrategias de Inicialización y Warm-Start}

La inicialización de parámetros variacionales determina en gran medida el éxito de la optimización VQE. La inicialización aleatoria conduce frecuentemente a mínimos locales y requiere cientos de iteraciones para converger, especialmente cerca del punto crítico donde el paisaje energético se vuelve complejo.

\subsection{Transferencia de Parámetros}

\citet{mele2022} demostraron teóricamente que los parámetros óptimos del HVA varían suavemente con los parámetros del Hamiltoniano. Esto fundamenta la estrategia de \emph{warm-start} por barrido descendente: optimizando primero en $h = 2$ (donde $|+\rangle^{\otimes N}$ es casi exacto y $\boldsymbol{\theta}^* \approx 0$) y propagando la solución hacia valores menores de $h$, cada nueva optimización parte de una región cercana al óptimo global.

\citet{puig2025} formalizaron este resultado en PRX Quantum, demostrando que el warm-start proporciona varianzas de pérdida mayores (señales de gradiente más fuertes) comparadas con la inicialización aleatoria. Cuantitativamente, la transferencia de parámetros reduce el número de iteraciones de optimización de $\sim 500$--$1000$ (inicialización aleatoria cerca del punto crítico) a $\sim 20$--$50$ iteraciones (warm-start desde el punto adyacente), constituyendo una aceleración de $10$--$50\times$ que se acumula a lo largo del barrido completo.

\citet{skogh2023} validaron experimentalmente esta estrategia para Hamiltonianos parametrizados, confirmando que la transferencia de parámetros entre instancias cercanas del Hamiltoniano preserva la convergencia al mínimo global.

\subsection{Predicción de Parámetros mediante Aprendizaje Automático}

La extensión natural del warm-start secuencial es la predicción directa de parámetros mediante modelos de aprendizaje automático. Tres trabajos recientes validan independientemente este paradigma:

\textbf{NN-VQE} \citep{miao2024}: Demostró que un perceptrón multicapa (MLP) entrenado con solo 20 puntos de datos puede predecir parámetros VQE óptimos para Hamiltonianos de espines parametrizados. Los autores emplean dropout para regularización y proponen un esquema de aprendizaje activo para identificar regiones de alta incertidumbre. Este trabajo valida directamente la viabilidad de nuestro enfoque de predicción neuronal.

\textbf{Qracle} \citep{zhang2025}: Propuso un inicializador basado en GNN que codifica conjuntamente el Hamiltoniano y el ansatz como un grafo unificado. El modelo reduce hasta un 64\% las iteraciones de optimización VQE y generaliza a sistemas no vistos durante el entrenamiento. Este trabajo valida específicamente la elección de GNN (frente a MLP) para la predicción de parámetros variacionales.

\textbf{Flow-VQE} \citep{zou2026}: Empleó flujos normalizantes generativos para el warm-start de VQE, alcanzando aceleraciones de hasta $50\times$. Aunque su enfoque generativo difiere de nuestro mapeo determinista, confirma que la predicción clásica de parámetros es una estrategia viable y competitiva.

Adicionalmente, \citet{lee2026} demostraron que un autoencoder de grafos combinado con redes neuronales puede generar parámetros VQE que generalizan entre instancias de Hamiltonianos sin optimización específica por instancia, validando directamente el paradigma MPNN-para-VQE que adoptamos.

\section{Redes Neuronales de Grafos}

\subsection{Fundamentos del Paso de Mensajes}

Las Redes Neuronales de Grafos (GNN) operan sobre datos estructurados como grafos $\mathcal{G} = (V, E)$, donde $V$ es el conjunto de nodos y $E$ el conjunto de aristas. El paradigma de Paso de Mensajes (Message Passing), formalizado por \citet{gilmer2017}, define la actualización de representaciones nodales mediante la agregación iterativa de información de vecinos:

\begin{equation}
\mathbf{h}_v^{(l+1)} = \text{UPDATE}\left(\mathbf{h}_v^{(l)}, \text{AGGREGATE}\left(\{\mathbf{h}_u^{(l)} : u \in \mathcal{N}(v)\}\right)\right)
\label{eq:mpnn}
\end{equation}

\noindent donde $\mathbf{h}_v^{(l)}$ es la representación del nodo $v$ en la capa $l$, $\mathcal{N}(v)$ denota los vecinos de $v$, y las funciones UPDATE y AGGREGATE son parametrizadas por redes neuronales.

Para sistemas de espines, la representación como grafo es natural e isomórfica: los qubits son nodos, las interacciones físicas ($J_{ij}$) son aristas, y los campos locales ($h_i$) son atributos de nodo. Esta correspondencia directa ---sin pérdida de información ni overhead de codificación--- es una ventaja fundamental de las GNN frente a arquitecturas alternativas como CNN o MLP para problemas de muchos cuerpos \citep{kochkov2021}.

\subsection{GINConv: Máxima Expresividad para Grafos}

\citet{xu2019} demostraron en un resultado teórico fundamental que la Graph Isomorphism Network (GIN) es tan poderosa como el test de Weisfeiler-Lehman (WL) para distinguir grafos no isomorfos. Esto la convierte en la arquitectura MPNN maximalmente expresiva. La actualización GINConv se define como:

\begin{equation}
\mathbf{h}_v^{(l+1)} = \text{MLP}^{(l)}\left((1 + \epsilon^{(l)}) \cdot \mathbf{h}_v^{(l)} + \sum_{u \in \mathcal{N}(v)} \mathbf{h}_u^{(l)}\right)
\label{eq:ginconv}
\end{equation}

\noindent donde $\epsilon^{(l)}$ es un parámetro aprendible que controla la importancia relativa del nodo central frente a sus vecinos.

Para redes uniformes (como la cadena 1D del TFIM con $J$ constante), todas las aristas son equivalentes y los mecanismos de atención (como GATConv) añaden parámetros sin ganancia informativa. \citet{meng2025} confirmaron independientemente que las GNN superan a las CNN en un 36\% para la predicción de propiedades de circuitos cuánticos, validando la superioridad de las arquitecturas basadas en grafos para este dominio.

\subsection{GNN para Sistemas de Espines}

La aplicación de GNN a sistemas de espines ha sido validada por múltiples trabajos recientes:

\citet{kochkov2021} emplearon GNN como manifold variacional para estados fundamentales de Hamiltonianos de Heisenberg diversos, demostrando que la arquitectura respeta simetrías físicas por construcción y generaliza a sistemas de mayor tamaño.

\citet{slavin2025} aplicaron GNN específicamente a la predicción de magnetización en sistemas de Ising cuasi-unidimensionales, capturando correctamente mesetas de magnetización, transiciones críticas y efectos de frustración geométrica a partir de la geometría de la red. Este trabajo valida directamente nuestro paradigma de predicción GNN-desde-grafo para sistemas de espines.

\citet{huang2022} proporcionaron la fundamentación teórica más rigurosa, demostrando que el aprendizaje automático clásico puede predecir eficientemente propiedades del estado fundamental de Hamiltonianos con gap, generalizando a nuevas instancias dentro de la misma fase. Este resultado establece la base teórica para la capacidad de generalización de nuestra MPNN a lo largo del barrido en $h$.

\section{Mitigación de Errores en Hardware NISQ}

La ejecución en hardware cuántico real introduce tres fuentes de degradación: errores de compuerta ($\sim 0.1$--$0.5\%$ por compuerta de dos qubits), ruido de muestreo (shot noise $\sim 1/\sqrt{N_{\text{shots}}}$) y errores de lectura ($\sim 1$--$3\%$ por qubit). Para que el pipeline produzca resultados físicamente significativos en hardware, es necesario emplear técnicas de mitigación de errores.

\subsection{Zero Noise Extrapolation (ZNE)}

La ZNE ejecuta el mismo circuito a múltiples niveles de ruido artificialmente amplificado (típicamente $1\times$, $2\times$, $3\times$) y extrapola al límite de ruido cero. \citet{uvarov2024} propusieron una variante inhomogénea que explota la distribución no uniforme de errores de compuerta en el hardware real, utilizando diferentes mappings de qubits para generar niveles de ruido distintos sin insertar compuertas adicionales. Esta técnica es directamente aplicable a procesadores IBM donde las tasas de error varían significativamente entre pares de qubits.

\citet{sun2025} extendieron la ZNE reemplazando la extrapolación polinomial clásica por una red neuronal entrenada para predecir el valor libre de ruido, reduciendo los errores residuales a $O(10^{-2})$. Este enfoque de ``ZNE mejorada por NN'' está implementado en nuestro pipeline como módulo \texttt{NNExtrapolator}.

\subsection{Dynamical Decoupling (DD)}

El Dynamical Decoupling inserta secuencias de pulsos de identidad durante los periodos ociosos del circuito para suprimir la decoherencia por acoplamiento con el entorno. \citet{pokharel2025} demostraron la optimización de secuencias DD mediante algoritmos genéticos en procesadores IBM de hasta 100 qubits, logrando supresión de errores escalable sin costo adicional en shots. Esta técnica es particularmente efectiva para nuestros circuitos HVA, donde los periodos ociosos entre capas de compuertas son significativos.

\subsection{Validación en Hardware Contemporáneo}

Varios trabajos recientes validan la viabilidad de VQE en hardware IBM actual:

\citet{kiiamov2026} ejecutaron VQE de 6 qubits en IBM Heron 2 para localización de Wigner, alcanzando errores relativos menores al 7\%. Este resultado confirma que el hardware actual puede producir resultados VQE significativos a nuestra escala de sistema ($N = 6$--$10$).

\citet{sharma2026} compararon diagonalización exacta, VQE simulado y ejecución en hardware IQM Garnet para el TFIM 1D, demostrando que el ruido ensancha la región de crossover crítico pero preserva la clasificación correcta de fases. Este resultado es directamente relevante para nuestra Fase 4: el éxito se mide por clasificación de fase correcta, no por fidelidad del estado.

A escala de utilidad, \citet{ahsan2025} ejecutaron VQE de 103 sitios en la red de Kagomé sobre IBM Heron, alcanzando energías por sitio de $-0.417J$ consistentes con el límite termodinámico. Aunque nuestro trabajo opera a escalas menores ($N = 6$--$10$), este resultado demuestra que el pipeline híbrido (pre-optimización clásica + ejecución cuántica superficial + mitigación) es viable a escala de utilidad.

\section{Sistemas de Espines vs. Química Cuántica: Justificación Estratégica}

Una decisión arquitectónica fundamental de este trabajo es la orientación hacia sistemas de espines de materia condensada, en lugar de la química cuántica molecular. Esta elección se justifica por la eficiencia de mapeo al hardware cuántico:

\begin{itemize}
    \item \textbf{Espines:} Existe un isomorfismo natural entre un espín-1/2 y un qubit. Las interacciones $Z_i Z_j$ se traducen directamente en compuertas nativas $R_{ZZ}(\theta)$. Para el TFIM con $N$ sitios y HVA $p = 2$: el circuito requiere exactamente $2(N-1)$ compuertas de dos qubits ---lineal en el tamaño del sistema.
    
    \item \textbf{Moléculas:} Los electrones son fermiones que obedecen relaciones de anticonmutación. La transformación de Jordan-Wigner mapea operadores fermiónicos locales a cadenas de Pauli de longitud $O(N)$, incrementando la profundidad del circuito a $O(N^4)$ para ansätze tipo UCCSD ---muy por encima del límite de truncamiento $O(\log n)$ de \citet{mele2026}.
\end{itemize}

Adicionalmente, la clasificación de fases cuánticas es una pregunta \emph{cualitativa} (¿en qué fase está el sistema?) con un criterio de éxito relajado ($\Delta E/\text{gap} < 5\%$), mientras que la química cuántica exige ``precisión química'' ($\Delta E < 1.6 \times 10^{-3}$ Hartree), inalcanzable con circuitos superficiales en hardware ruidoso.

\section{Conclusiones del Capítulo}

El análisis del estado del arte revela una convergencia teórica que habilita el enfoque propuesto en este trabajo:

\begin{enumerate}
    \item El truncamiento por ruido no unital \citep{mele2026} impone circuitos superficiales como \emph{requisito}, no como limitación. El HVA $p \leq 2$ opera dentro de este régimen por diseño.
    
    \item Las funciones de costo locales garantizan gradientes detectables en circuitos superficiales \citep{cerezo2021}, eliminando los barren plateaus como obstáculo.
    
    \item La transferibilidad de parámetros en el HVA \citep{mele2022, puig2025} y la capacidad demostrada de las GNN para predecir propiedades de estados fundamentales \citep{huang2022, zhang2025} fundamentan la viabilidad del predictor MPNN.
    
    \item Los sistemas de espines ofrecen un mapeo isomórfico a qubits que mantiene los circuitos dentro del régimen coherente del hardware actual, a diferencia de la química cuántica.
    
    \item La mitigación de errores (ZNE + DD + TREX) ha sido validada experimentalmente en hardware IBM a escalas comparables a nuestro sistema \citep{kiiamov2026, sharma2026}.
\end{enumerate}

La brecha identificada en la literatura es la ausencia de un pipeline \emph{integrado} que combine todas estas técnicas en una arquitectura end-to-end: desde la generación de datos clásicos hasta el despliegue en hardware con clasificación de fases validada. Los trabajos existentes abordan componentes individuales (predicción de parámetros, mitigación de errores, diseño de ansatz), pero ninguno demuestra la cadena completa con validación cuantitativa en múltiples escalas de sistema. Este es precisamente el vacío que el presente trabajo pretende llenar.


%---------------------------
% CAPÍTULO 3: OBJETIVOS Y METODOLOGÍA DE TRABAJO
%---------------------------

\chapter{Objetivos y Metodología de Trabajo}

Este capítulo establece el puente entre el estudio del dominio presentado en el capítulo anterior y la contribución específica de este trabajo. Se formulan el objetivo general y los objetivos específicos siguiendo el criterio SMART \citep{doran1981}, y se describe la metodología de trabajo adoptada para alcanzarlos.

\section{Objetivo General}

Demostrar que un pipeline híbrido clásico-cuántico ---basado en la predicción de parámetros variacionales mediante una Red Neuronal de Grafos de Paso de Mensajes (MPNN) inyectados en un Ansatz Variacional Hamiltoniano (HVA) superficial--- permite clasificar correctamente fases cuánticas del Modelo de Ising en Campo Transversal con un error relativo $\Delta E/\text{gap} < 5\%$ y un costo cuántico de 0--2 iteraciones de refinamiento, validado en simulación sin ruido para cadenas de $N = 6$ y $N = 10$ espines, y desplegable en hardware IBM con mitigación de errores.

Este objetivo es:
\begin{itemize}
    \item \textbf{Específico:} Define el sistema (TFIM), la métrica ($\Delta E/\text{gap} < 5\%$), la escala ($N = 6$, $N = 10$) y la plataforma (IBM).
    \item \textbf{Medible:} El criterio de éxito es cuantitativo: checklist de 4 métricas con umbrales definidos.
    \item \textbf{Alcanzable:} Basado en resultados preliminares que demuestran viabilidad (40+ experimentos a $N = 6$).
    \item \textbf{Relevante:} Aborda la brecha identificada en el estado del arte (ausencia de pipeline integrado end-to-end).
    \item \textbf{Temporal:} Acotado al periodo del TFM (semestre académico 2025--2026).
\end{itemize}

\section{Objetivos Específicos}

\begin{enumerate}
    \item \textbf{OE1 --- Analizar} el estado del arte en algoritmos cuánticos variacionales, identificando las barreras teóricas (barren plateaus, truncamiento por ruido) y las estrategias de inicialización inteligente propuestas en la literatura reciente (2022--2026), con énfasis en enfoques basados en redes neuronales de grafos.
    
    \item \textbf{OE2 --- Desarrollar} un módulo de generación de ground truth clásico que produzca estados fundamentales exactos, energías y observables locales del TFIM 1D para $N = 6$ y $N = 10$ espines mediante diagonalización exacta y DMRG, con un barrido no uniforme del campo transversal $h \in [0, 2]$ (mayor densidad en la región crítica $h \in [0.8, 1.4]$).
    
    \item \textbf{OE3 --- Implementar} un pipeline de optimización VQE con warm-start descendente ($h = 2 \to 0$) sobre un HVA de profundidad $p = 2$, alcanzando fidelidades $\geq 93\%$ respecto al estado fundamental exacto en al menos el 90\% de los puntos del barrido, con un máximo de 5 reinicios por punto.
    
    \item \textbf{OE4 --- Diseñar y entrenar} un predictor MPNN basado en GINConv que, dado un grafo representando el Hamiltoniano del TFIM, prediga los parámetros óptimos $\theta^*$ del HVA con un error cuadrático medio (MSE) inferior a $5 \times 10^{-3}$ sobre el conjunto de validación, verificando que la energía resultante satisfaga $\Delta E/\text{gap} < 5\%$ para $h \geq 1.25$ ($N = 6$) y $h \geq 1.4$ ($N = 10$).
    
    \item \textbf{OE5 --- Validar} el pipeline completo en simulación sin ruido para $N = 6$ (mínimo 40 experimentos, 14 configuraciones) y $N = 10$ (mínimo 15 experimentos, 3 semillas), documentando la reproducibilidad estadística y los límites físicos del enfoque (frontera de expresividad del HVA $p = 2$).
    
    \item \textbf{OE6 --- Implementar} el módulo de despliegue en hardware con mitigación de errores (ZNE inhomogéneo, Dynamical Decoupling, TREX) y verificar que el criterio $\Delta E/\text{gap} < 5\%$ con clasificación correcta de fase se mantiene bajo condiciones de ruido realistas, utilizando simulación ruidosa y/o ejecución en procesadores IBM Torino/Heron.
\end{enumerate}

\section{Metodología de Trabajo}

La metodología adoptada sigue un enfoque experimental iterativo, organizado en cuatro fases de desarrollo que se corresponden con las fases del pipeline propuesto. Cada fase produce artefactos verificables que alimentan la siguiente, permitiendo validación incremental.

\subsection{Fase 1: Generación del Ground Truth Clásico}

\textbf{Objetivo:} Producir el conjunto de datos de referencia que servirá como base de entrenamiento y validación para todo el pipeline.

\textbf{Instrumentos:}
\begin{itemize}
    \item Diagonalización exacta mediante \texttt{numpy.linalg.eigh} para $N \leq 10$ (espacio de Hilbert $\leq 1024$ dimensiones).
    \item DMRG mediante la librería TeNPy \citep{hauschild2018} como ruta alternativa para validación cruzada y escalabilidad futura.
    \item Construcción del Hamiltoniano como \texttt{SparsePauliOp} (Qiskit 2.x) para compatibilidad directa con las fases posteriores.
\end{itemize}

\textbf{Procedimiento:}
\begin{enumerate}
    \item Definir la topología de la red (cadena 1D, $N$ sitios, condiciones de contorno abiertas).
    \item Construir el Hamiltoniano TFIM (Ecuación~\ref{eq:tfim}) para un barrido no uniforme de $h$: paso $\Delta h = 0.1$ en las regiones $h \in [0, 0.8] \cup [1.4, 2.0]$ y paso $\Delta h = 0.05$ en la región crítica $h \in [0.8, 1.4]$, totalizando 27 puntos.
    \item Para cada $h$: calcular el estado fundamental $|\psi_0\rangle$, la energía $E_0$, el gap espectral $\Delta = E_1 - E_0$, y los observables locales $\langle X_i \rangle$, $\langle Z_i Z_{i+1} \rangle$.
    \item Almacenar el dataset con metadatos de integridad (versión, función de costo, checksums).
\end{enumerate}

\textbf{Criterio de éxito:} Dataset completo con 27 puntos, energías verificadas contra solución analítica (error $< 10^{-12}$), gap espectral positivo en todos los puntos.

\subsection{Fase 2: Optimización VQE con Warm-Start}

\textbf{Objetivo:} Obtener los parámetros óptimos $\theta^*(h)$ del HVA para cada punto del barrido, generando el dataset de entrenamiento para la MPNN.

\textbf{Instrumentos:}
\begin{itemize}
    \item Circuito HVA de $p = 2$ capas construido con compuertas $R_{ZZ}(2\theta)$ y $R_X(2\theta)$.
    \item Evaluación de energía mediante \texttt{StatevectorEstimator} (Qiskit Primitives V2) en simulación sin ruido.
    \item Optimizador L-BFGS-B con tolerancia $f_{\text{tol}} = 10^{-14}$ y límites $\theta \in [-\pi, \pi]$.
    \item Estrategia multi-restart: 5 reinicios con perturbación gaussiana $\sigma = 0.1$ alrededor del warm-start.
\end{itemize}

\textbf{Procedimiento:}
\begin{enumerate}
    \item Inicializar $\theta_0 = \text{uniform}(-0.01, 0.01)$ para $h = 2.0$ (estado inicial $|+\rangle^{\otimes N}$ es casi exacto).
    \item Para cada $h$ en orden descendente ($h = 2.0, 1.9, \ldots, 0.0$):
    \begin{enumerate}
        \item Heredar $\theta^*_{h+\Delta h}$ como semilla (warm-start).
        \item Ejecutar 5 optimizaciones independientes (semilla + 4 perturbaciones aleatorias).
        \item Seleccionar la solución de menor energía.
        \item Calcular fidelidad $F = |\langle \psi_0 | \psi(\theta^*) \rangle|^2$ contra el ground truth de Fase 1.
    \end{enumerate}
    \item Filtrar puntos con fidelidad $< 0.93$ (excluidos del entrenamiento MPNN).
\end{enumerate}

\textbf{Criterio de éxito:} Fidelidad media $\geq 0.95$ sobre el barrido completo; $\geq 90\%$ de puntos con fidelidad $\geq 0.93$; paisaje $\theta(h)$ suave (sin discontinuidades).

\subsection{Fase 3: Entrenamiento del Predictor MPNN}

\textbf{Objetivo:} Entrenar una red neuronal de grafos que prediga $\theta^*$ directamente desde la representación en grafo del Hamiltoniano, eliminando la necesidad de optimización VQE en tiempo de inferencia.

\textbf{Instrumentos:}
\begin{itemize}
    \item Arquitectura MPNN: 3 capas GINConv \citep{xu2019} + global\_mean\_pool + MLP de salida.
    \item Framework: PyTorch Geometric \citep{fey2019}.
    \item Representación del grafo: nodos = qubits (atributo: $h_i$), aristas = interacciones (atributo: $J_{ij}$).
    \item Hiperparámetros: hidden\_dim $= 64$ ($N = 6$) / $128$ ($N = 10$), epochs $= 6000$, lr $= 10^{-3}$, dropout $= 0.1$, patience $= 300$/$500$.
\end{itemize}

\textbf{Procedimiento:}
\begin{enumerate}
    \item Construir el dataset de grafos a partir de los pares $(h, \theta^*(h))$ filtrados por fidelidad $\geq 0.93$.
    \item Dividir en entrenamiento (80\%) y validación (20\%) con estratificación por fase.
    \item Entrenar minimizando MSE entre $\theta_{\text{pred}}$ y $\theta^*$.
    \item Validación física: para cada predicción, evaluar la energía $E(\theta_{\text{pred}})$ y verificar $\Delta E/\text{gap} < 5\%$.
    \item Análisis de gradientes de pesos \citep{hernandes2025}: detectar picos en $\|\partial W / \partial h\|$ como indicador de transición de fase codificada en la red.
\end{enumerate}

\textbf{Criterio de éxito:} MSE $< 5 \times 10^{-3}$; $\Delta E/\text{gap} < 5\%$ para $h \geq 1.25$ ($N = 6$) y $h \geq 1.4$ ($N = 10$); detección de pico de gradiente en $h \approx 1.0$--$1.2$.

\subsection{Fase 4: Despliegue en Hardware con Mitigación de Errores}

\textbf{Objetivo:} Demostrar que los parámetros predichos por la MPNN producen clasificación correcta de fases cuando se ejecutan en hardware cuántico real (o simulación ruidosa realista).

\textbf{Instrumentos:}
\begin{itemize}
    \item \texttt{EstimatorV2} de Qiskit IBM Runtime para ejecución en hardware.
    \item ZNE inhomogéneo: múltiples layouts de qubits con diferentes tasas de error efectivas \citep{uvarov2024}.
    \item Dynamical Decoupling: secuencias optimizadas insertadas mediante pass manager de Qiskit \citep{pokharel2025}.
    \item TREX: aleatorización de errores de lectura para corrección estadística.
    \item Extrapolador neuronal (\texttt{NNExtrapolator}): MLP entrenado para extrapolar a ruido cero \citep{sun2025}.
\end{itemize}

\textbf{Procedimiento:}
\begin{enumerate}
    \item Seleccionar puntos de test representativos: $h = 1.25$, $1.4$, $1.5$ (región de crossover y fase paramagnética).
    \item Para cada punto de test:
    \begin{enumerate}
        \item Obtener $\theta_{\text{pred}}$ de la MPNN entrenada.
        \item Construir circuito HVA con parámetros predichos.
        \item Transpilar al backend objetivo con \texttt{generate\_preset\_pass\_manager(optimization\_level=2)}.
        \item Ejecutar con ZNE (niveles de ruido $1\times$, $2\times$, $3\times$) + DD + TREX.
        \item Extrapolar observables a ruido cero.
        \item Evaluar checklist de 4 métricas.
    \end{enumerate}
    \item Opcionalmente: ejecutar 0--2 iteraciones de refinamiento adaptativo si $\Delta E/\text{gap} > 5\%$.
\end{enumerate}

\textbf{Criterio de éxito (checklist de 4 métricas):}
\begin{enumerate}
    \item $\Delta E/\text{gap} < 5\%$ (métrica primaria).
    \item Errores en observables locales $|\langle X_i \rangle - \langle X_i \rangle_{\text{exact}}| < 10^{-2}$.
    \item Clasificación correcta de fase (paramagnética vs. ferromagnética).
    \item Circuito dentro del límite de profundidad ($p \leq 2$, sin iteraciones ADAPT adicionales).
\end{enumerate}

\subsection{Análisis de Resultados}

Los resultados se analizan mediante:

\begin{itemize}
    \item \textbf{Estadística descriptiva:} Media, desviación estándar y distribución de métricas sobre múltiples semillas aleatorias (mínimo 3 semillas por configuración).
    \item \textbf{Reproducibilidad:} Cada experimento se ejecuta con semillas fijas y se documenta en bitácoras estructuradas (binnacles) con metadatos completos (versión de código, paquetes, plataforma).
    \item \textbf{Comparación con baselines:} VQE sin warm-start (inicialización aleatoria), VQE con warm-start pero sin MPNN, y predicción MLP (sin estructura de grafo).
    \item \textbf{Análisis de límites físicos:} Identificación de la frontera de expresividad del HVA $p = 2$ mediante barrido sistemático de $h$ y comparación con resultados de \citet{tripathi2026}.
    \item \textbf{Escalabilidad:} Comparación de métricas entre $N = 6$ y $N = 10$ para verificar que el pipeline no degrada con el tamaño del sistema.
\end{itemize}

\subsection{Herramientas y Entorno de Desarrollo}

\begin{itemize}
    \item \textbf{Lenguaje:} Python 3.11+.
    \item \textbf{Computación cuántica:} Qiskit 2.x (Primitives V2, \texttt{SparsePauliOp}, pass managers).
    \item \textbf{Hardware cuántico:} IBM Quantum Platform (procesadores Torino/Heron, 133+ qubits).
    \item \textbf{Redes tensoriales:} TeNPy \citep{hauschild2018} para DMRG.
    \item \textbf{Machine Learning:} PyTorch + PyTorch Geometric \citep{fey2019} para la MPNN.
    \item \textbf{Optimización:} SciPy (L-BFGS-B) para VQE clásico.
    \item \textbf{Gestión de experimentos:} Scripts automatizados con registro estructurado (JSON), bitácoras markdown, y pruebas de humo (\texttt{smoke\_test}) para validación continua.
    \item \textbf{Control de calidad:} Suite de 33 tests unitarios (pytest), linter (ruff), hooks pre-commit para verificación de restricciones arquitectónicas (profundidad HVA, primitivas V2).
\end{itemize}


\begin{thebibliography}{99}

\bibitem[Ahsan et al.(2025)]{ahsan2025} Ahsan, M. et al. (2025). \emph{Utility-scale quantum computation of ground-state energy in a 100+ site planar Kagome antiferromagnet via Hamiltonian engineering}. arXiv preprint arXiv:2507.06361.

\bibitem[Anderson(1972)]{anderson1972} Anderson, P. W. (1972). \emph{More is different}. Science, 177(4047), 393-396.

\bibitem[Anderson(1973)]{anderson1973} Anderson, P. W. (1973). \emph{Resonating valence bonds: A new kind of insulator?}. Materials Research Bulletin, 8(2), 153-160.

\bibitem[Cerezo et al.(2021)]{cerezo2021} Cerezo, M., Sone, A., Volkoff, T., Cincio, L., \& Coles, P. J. (2021). \emph{Cost function dependent barren plateaus in shallow parametrized quantum circuits}. Nature Communications, 12(1), 1791.

\bibitem[Doran(1981)]{doran1981} Doran, G. T. (1981). \emph{There's a S.M.A.R.T. way to write management's goals and objectives}. Management Review, 70(11), 35-36.

\bibitem[Dutta et al.(2015)]{dutta2015} Dutta, A., Aeppli, G., Chakrabarti, B. K., Divakaran, U., Rosenbaum, T. F., \& Sen, D. (2015). \emph{Quantum phase transitions in transverse field spin models: From statistical physics to quantum information}. Cambridge University Press.

\bibitem[Feynman(1982)]{feynman1982} Feynman, R. P. (1982). \emph{Simulating physics with computers}. International Journal of Theoretical Physics, 21(6/7), 467-488.

\bibitem[Fey \& Lenssen(2019)]{fey2019} Fey, M., \& Lenssen, J. E. (2019). \emph{Fast graph representation learning with PyTorch Geometric}. ICLR Workshop on Representation Learning on Graphs and Manifolds.

\bibitem[Gilmer et al.(2017)]{gilmer2017} Gilmer, J., Schoenholz, S. S., Riley, P. F., Vinyals, O., \& Dahl, G. E. (2017). \emph{Neural message passing for quantum chemistry}. Proceedings of the 34th International Conference on Machine Learning (ICML), 70, 1263-1272.

\bibitem[Hauschild \& Pollmann(2018)]{hauschild2018} Hauschild, J., \& Pollmann, F. (2018). \emph{Efficient numerical simulations with Tensor Networks: Tensor Network Python (TeNPy)}. SciPost Physics Lecture Notes, 5.

\bibitem[Hernandes et al.(2025)]{hernandes2025} Hernandes, V. et al. (2025). \emph{Adiabatic fine-tuning of neural quantum states enables detection of phase transitions in weight space}. arXiv preprint arXiv:2503.17140v2.

\bibitem[Huang et al.(2022)]{huang2022} Huang, H.-Y., Kueng, R., Torlai, G., Albert, V. V., \& Preskill, J. (2022). \emph{Provably efficient machine learning for quantum many-body problems}. Science, 377(6613), eabk3333.

\bibitem[Kiiamov et al.(2026)]{kiiamov2026} Kiiamov, A. G. et al. (2026). \emph{Simulating Wigner localisation with the IBM Heron 2 quantum processor: A proof-of-principle benchmarking study}. arXiv preprint arXiv:2601.01263.

\bibitem[Kochkov et al.(2021)]{kochkov2021} Kochkov, D., Pfaff, T., Sanchez-Gonzalez, A., Battaglia, P., \& Clark, B. K. (2021). \emph{Learning ground states of quantum Hamiltonians with graph networks}. arXiv preprint arXiv:2110.06390.

\bibitem[Lee et al.(2026)]{lee2026} Lee, J. et al. (2026). \emph{Improving generalization and trainability of quantum eigensolvers via graph neural encoding}. arXiv preprint arXiv:2602.19752.

\bibitem[McClean et al.(2018)]{mcclean2018} McClean, J. R., Boixo, S., Smelyanskiy, V. N., Babbush, R., \& Neven, H. (2018). \emph{Barren plateaus in quantum neural network training landscapes}. Nature Communications, 9(1), 4812.

\bibitem[Mele et al.(2022)]{mele2022} Mele, A. A., Mbeng, G. B., Santoro, G. E., Collura, M., \& Torta, P. (2022). \emph{Avoiding barren plateaus via transferability of smooth solutions in a Hamiltonian variational ansatz}. Physical Review A, 106, L060401.

\bibitem[Mele et al.(2026)]{mele2026} Mele, A. A., Angrisani, A., Ghosh, S., Khatri, S., Eisert, J., Stilck França, D., \& Quek, Y. (2026). \emph{Noise-induced shallow circuits and the absence of barren plateaus}. Nature Physics.

\bibitem[Meng et al.(2025)]{meng2025} Meng, F. et al. (2025). \emph{Output prediction of quantum circuits based on graph neural networks}. arXiv preprint arXiv:2504.00464.

\bibitem[Miao et al.(2024)]{miao2024} Miao, J., Hsieh, C.-Y., \& Zhang, S.-X. (2024). \emph{Neural-network-encoded variational quantum algorithms}. Physical Review Applied, 21, 014053.

\bibitem[Peruzzo et al.(2014)]{peruzzo2014} Peruzzo, A., McClean, J., Shadbolt, P., Yung, M.-H., Zhou, X.-Q., Love, P. J., Aspuru-Guzik, A., \& O'Brien, J. L. (2014). \emph{A variational eigenvalue solver on a photonic quantum processor}. Nature Communications, 5, 4213.

\bibitem[Pokharel et al.(2025)]{pokharel2025} Pokharel, B. et al. (2025). \emph{Empirical learning of dynamical decoupling on quantum processors}. arXiv preprint arXiv:2403.02294v2.

\bibitem[Preskill(2018)]{preskill2018} Preskill, J. (2018). \emph{Quantum computing in the NISQ era and beyond}. Quantum, 2, 79.

\bibitem[Puig et al.(2025)]{puig2025} Puig, R., Drudis, M., Thanasilp, S., \& Holmes, Z. (2025). \emph{Variational quantum simulation: A case study for understanding warm starts}. PRX Quantum, 6, 010317.

\bibitem[Savary \& Balents(2016)]{savary2016} Savary, L., \& Balents, L. (2016). \emph{Quantum spin liquids: a review}. Reports on Progress in Physics, 80(1), 016502.

\bibitem[Schollwöck(2011)]{schollwock2011} Schollwöck, U. (2011). \emph{The density-matrix renormalization group in the age of matrix product states}. Annals of Physics, 326(1), 96-192.

\bibitem[Sharma(2026)]{sharma2026} Sharma, R. (2026). \emph{Quantum phase transitions in the transverse-field Ising model: A comparative study of exact, variational, and hardware-based approaches}. arXiv preprint arXiv:2601.17515.

\bibitem[Skogh et al.(2023)]{skogh2023} Skogh, M., Leinonen, O., Lolur, P., \& Rahm, M. (2023). \emph{Accelerating variational quantum eigensolver convergence using parameter transfer}. Electronic Structure, 5, 035002.

\bibitem[Slavin(2025)]{slavin2025} Slavin, V. (2025). \emph{Graph neural network approach to predicting magnetization in quasi-one-dimensional Ising systems}. arXiv preprint arXiv:2507.17509.

\bibitem[Sun et al.(2025)]{sun2025} Sun, W. et al. (2025). \emph{Noise-mitigated variational quantum eigensolver with pre-training and zero-noise extrapolation}. arXiv preprint arXiv:2501.01646.

\bibitem[Tilly et al.(2022)]{tilly2022} Tilly, J., Chen, H., Cao, S., Picozzi, D., Setia, K., Li, Y., Grant, E., Wossnig, L., Rungger, I., Booth, G. H., \& Tennyson, J. (2022). \emph{The variational quantum eigensolver: A review of methods and best practices}. Physics Reports, 986, 1-128.

\bibitem[Tripathi et al.(2026)]{tripathi2026} Tripathi, A. P., Mathur, N., \& Tripathi, V. (2026). \emph{Ansätz expressivity and optimization in variational quantum simulations of transverse-field Ising model across system sizes}. arXiv preprint arXiv:2604.20961.

\bibitem[Troyer \& Wiese(2005)]{troyer2005} Troyer, M., \& Wiese, U. J. (2005). \emph{Computational complexity and fundamental limitations to fermionic quantum Monte Carlo simulations}. Physical Review Letters, 94(17), 170201.

\bibitem[Uvarov et al.(2024)]{uvarov2024} Uvarov, A. et al. (2024). \emph{Mitigating quantum gate errors for variational eigensolvers using hardware-inspired zero-noise extrapolation}. arXiv preprint arXiv:2307.11156v3.

\bibitem[Wiersema et al.(2020)]{wiersema2020} Wiersema, R., Zhou, C., de Sereville, Y., Carrasquilla, J., \& Kim, Y. B. (2020). \emph{Exploring entanglement and optimization within the Hamiltonian Variational Ansatz}. PRX Quantum, 1(2), 020319.

\bibitem[Xu et al.(2019)]{xu2019} Xu, K., Hu, W., Leskovec, J., \& Jegelka, S. (2019). \emph{How powerful are graph neural networks?}. Proceedings of the 7th International Conference on Learning Representations (ICLR).

\bibitem[Zhang et al.(2025)]{zhang2025} Zhang, C., Jiang, L., \& Chen, F. (2025). \emph{Qracle: A graph-neural-network-based parameter initializer for variational quantum eigensolvers}. arXiv preprint arXiv:2505.01236.

\bibitem[Zou et al.(2026)]{zou2026} Zou, H., Rahm, M., Kockum, A. F., \& Olsson, S. (2026). \emph{Generative flow-based warm start of the variational quantum eigensolver}. npj Quantum Information, 12, 5.

\end{thebibliography}

\end{document}
